\Askhsh{20672_SOLUTION}{1}
\begin{ylh}{1}
    \item [Ύλη από εικόνα]
\end{ylh}
ΛΥΣΗ
\begin{alist}
    \item [α)] Η ακτίνα του ημικυκλίου $C_1$ είναι $R_1 = \frac{AB}{2} = \frac{6}{2} = 3$ και το εμβαδόν του ισούται με
    \[ E_1 = \frac{\pi R_1^2}{2} = \frac{9\pi}{2}. \]
    Η ακτίνα του ημικυκλίου $C_2$ είναι $R_2 = \frac{B\Gamma}{2} = \frac{4}{2} = 2$ και το εμβαδόν του ισούται με
    \[ E_2 = \frac{\pi R_2^2}{2} = \frac{4\pi}{2} = 2\pi. \]
    Η ακτίνα του ημικυκλίου $C_3$ είναι $R_3 = \frac{A\Gamma}{2} = \frac{2}{2} = 1$ και το εμβαδόν του ισούται με
    \[ E_3 = \frac{\pi R_3^2}{2} = \frac{\pi}{2}. \]
    \item [β)] Το εμβαδόν του σκιασμένου χωρίου ισούται με
    \[ E_1 - E_2 + E_3 = \frac{9\pi}{2} - \frac{4\pi}{2} + \frac{\pi}{2} = \frac{6\pi}{2} = 3\pi. \]
\end{alist}